%Type of systems suitable for SA.
The degree to which a software must satisfy its requirements depends on the criticality of the software. An undetected fault in critical software may result in significant financial losses or human fatalities, e.g., a fault in a flight control system. Thus, the degree of the conformance, i.e., an adequacy criterion,  for critical software is set at the highest level -- software must comply with its requirements on all possible inputs. Yet, for the majority of applications adequacy criteria are more relaxed, e.g., to be $85\%$ statement adequate, a property must hold on a set of inputs that executes $85\%$ of the application's statements. 

To accommodate these different levels of compliance, software engineers have developed various techniques best suitable for a specific criterion. Mainly these techniques are classified into dynamic and static analyses. Dynamic analysis, such as testing, verifies that a program is adequate by executing it. This approach is widely used for non-critical software that uses weaker adequacy criteria. Employing testing to check a program's conformance on all inputs is usually intractable. Static analysis checks the conformance of a program to its property by analyzing the program code, i.e., without executing it. This compile-time determination allows static analysis to consider all possible program behaviors, which is suitable for checking program conformance to its properties at the highest level. Hence, static analysis is a key technique for verifying critical systems.
 
%Not enough precise SA. False positives.
The power of static analysis to consider all program behaviors follows from its ability to safely approximate program behaviors by abstracting the concrete domain of program variables and the programming language semantics. But at the same time because of its over-approximating nature, static analysis may identify some property violations as uncertain. The reason for this uncertainty is that a static analysis cannot tell if a violation happens on a feasible or an infeasible program behavior. This inconclusiveness is unacceptable for critical system and thus each potential violation must be examined further. An automatic solution to the elimination of false positive violations is to increase the precision of a static analysis, i.e., improve the analysis so it considers fewer infeasible behaviors.

%SA can be make arbitrary precise by either making its domain finer (C\&C 77) or combining different analyses together (C\&C 79). Scalability issues.
%Scalability and precision as two opposing properties of SA. Ideally good scalability and good precision together. Two approach to maintain both.
One way to improve the precision of a static analysis is to refine its abstract domain and program semantics interpretation which is always theoretically possible~\cite{Cousot:1977}, or to combine different static analyses in a suitable manner as discussed in~\cite{Cousot:1979}. Besides the laborious effort of redesigning a static analysis or combining two of them together, the resulting analysis usually becomes expensive in terms of time and memory consumption which makes it inapplicable to large applications. Thus precision and scalability of a static analysis work against each other. When a static analysis has a high precision generally it is not scalable and conversely when a static analysis is scalable its precision typically suffers. The literature suggests two complementary ideas for approaching this dilemma.\\
One approach focuses on making a precise static analysis scalable by partitioning a large program into modules like procedures and classes, and allowing the analysis to examine each partition independently. Next, the analyzed information of each module is composed together to obtain the result of  the whole program analysis. In the literature~\cite{Cousot:2002, Ballabriga:2007} this method is referred to as {\em partial static analysis}.\\
%Can put Corina's there if it said that it can only handle couple modules
%Conditional analyses. 3 papers. Assumes some conditions holds and another analyses should verify that these conditions holds
Another approach concentrates on improving the precision of a scalable analysis by permitting the analysis to explore only those program states for which it is adequately precise, i.e., able to provide definitive result. In the literature~\cite{Naik:POPL07,Conway:SAS08, Beyer:Tech11} this approach is called {\em conditional static analysis} since the permitted states are described by a condition $\theta$. In this framework an analysis verifies a program under some assumptions, i.e., there is no null pointer exception in the program. Next, another analysis attempts to prove these assumptions by either showing that the states which do not satisfy $\theta$ are not reachable or do not lead to property violations. In the previous work the condition $\theta$ either determined from the analysis design~\cite{Naik:POPL07,Conway:SAS08} when $\theta$ is applicable to all program states or determined during program analysis execution~\cite{Beyer:Tech11} when $\theta$ is composed from the conditions assumed to be held for a certain set of states.\\
%State briefly what has been done in the previous work of conditional model checker and conditional soundness and what we are doing different there. Write related work first?
%Our work explores conditional data-flow analyses. What challenges or what has been accomplished.
Our work focuses on the latter approach applied to data-flow analysis, i.e., conditional data-flow analysis. Unlike the cases presented in the previous work, data-flow analysis cannot determine a priori or during its execution the conditions under which it is able to verify a program to its best abilities. Moreover, data-flow analysis has two alternative ways for expressing $\theta$. The main goal of our work is to explore different approaches in $\theta$ encoding for data-flow analysis so that a scalable analysis is able to optimize its result. In the next section we illustrate the concept of conditional data-flow analysis, different encodings for $\theta$ and their effect. In the next section we also present our research questions and the outline for the rest of the paper.
%The main difference is that condition is unknown a priori . Condition can be specified in multiple ways.